% Bewertungspaket zum Gesamttest «Flächen, Pythagoras und Ähnlichkeit» — Quelle des PDFs (python3 scripts/build-lp-pdf.py pythagoras).
% Ganze Punkte je Teilschritt; Werte und Fehlerwerte mit python3 nachgerechnet (scripts/lp/pythagoras/zahlen.py).
% Gerüst aus dem Bewertungspaket von LP 3 (gleichungen-formeln): Auftrag an die KI, Folgefehler, (E)-Zeilen.
\documentclass[11pt]{article}
\usepackage{../lp-druck}
\renewcommand{\lpname}{Flächen, Pythagoras und Ähnlichkeit}
\renewcommand{\lpteilgebiet}{}
\renewcommand{\lpfuss}{Bewertungspaket}
\newcolumntype{P}{>{\centering\arraybackslash}p{10mm}}
\newenvironment{raster}{\par\smallskip\noindent\renewcommand{\arraystretch}{1.6}\linespread{1.35}\selectfont\tabularx{\linewidth}{@{}>{\raggedright\arraybackslash}p{38mm}P X@{}}\toprule
  \small\textsc{Teilschritt} & \small\textsc{P} & \small\textsc{Musterlösung}\\\midrule}{\bottomrule\endtabularx\par}
\newcommand{\fehler}[1]{\par\smallskip{\small\textcolor{lprot}{\bfseries Typische Fehler:} #1}}
\newcommand{\E}{\,{\footnotesize\bfseries\color{lpgruen}(E)}}
\begin{document}
\lpkopf{Bewertungspaket zum Gesamttest}{}

\begin{lpachtung}{Erst lösen, dann öffnen}
Dieses Paket enthält die Musterlösung. Wer es vor dem Lösen liest, misst am Ende nur, wie gut das Abschreiben gelingt.
\end{lpachtung}

\section*{1 · So gehst du vor}
\begin{enumerate}[leftmargin=*,itemsep=2pt]
\item \textbf{Gesamttest lösen} — vollständig, mit Rechenweg, ohne dieses Paket und ohne Taschenrechner, mit Lineal.
\item \textbf{Selbst bewerten} mit Abschnitt~3: Zeile für Zeile vergleichen und die Punkte eintragen. Das geht ohne Hilfsmittel
  und ist der Normalfall.
\item \textbf{Oder bewerten lassen.} Ob du dafür eine KI nutzt und welche, entscheidest du selbst und in eigener Verantwortung —
  je nachdem, was dir aufgrund deines Alters und deiner persönlichen Situation erlaubt ist (Altersgrenzen und
  Nutzungsbedingungen des Anbieters, allenfalls Einverständnis deiner Eltern). Fotografiere oder scanne deine Lösung gut lesbar,
  \textbf{ohne Namen}. Handyfotos tragen oft unsichtbar Standort, Zeit und Gerät mit (Metadaten) — lade lieber einen Scan oder ein
  Bildschirmfoto des Fotos hoch. Lade nur deine Lösung, das Testblatt und dieses PDF hoch, keine persönlichen Daten, und füge den Auftrag aus
  Abschnitt~2 ein.
\item \textbf{Rückmeldung prüfen.} Stimmt, was die KI aus deiner Handschrift gelesen hat? Eine KI kann sich irren —
  im Zweifel gilt das Raster in Abschnitt~3.
\item \textbf{Gezielt wiederholen} nach Abschnitt~4.
\end{enumerate}

\section*{2 · Auftrag an die KI}
\begin{tcolorbox}[colback=lplinie!25,colframe=lplinie,boxrule=0.5pt,arc=1mm,fontupper=\small]
Du bist Korrektorin oder Korrektor für Mathematik auf der Sekundarstufe I in der Schweiz. Im Anhang findest du (1)~das Testblatt mit den Aufgaben, (2)~das Bewertungspaket zum Gesamttest «Flächen, Pythagoras und Ähnlichkeit» mit Musterlösung und Punkteraster und (3)~Bilder meiner handschriftlichen Lösung.\par\medskip
Geh so vor:\par
1. Schreib zu jeder Aufgabe G1 bis G8 zuerst ab, was du in meiner Lösung liest — Rechenweg, Ergebnis mit Einheit; bei Skizzen (G5, G6, G8), welches Dreieck eingezeichnet ist und wo der rechte Winkel liegt. Was du nicht sicher lesen kannst, markierst du mit [unsicher]. Nicht raten. Achte besonders auf Hochzahlen ($^2$), Wurzelzeichen, Dezimalpunkte und Einheiten (cm, m, cm$^2$).\par
2. Vergib die Punkte genau nach dem Raster im Bewertungspaket (Abschnitt 3), ganze Punkte je Zeile. Ziehe einen Fehler nur einmal ab: Rechne ich mit einem falschen Zwischenergebnis richtig weiter, gibt es die Folgepunkte. Ausnahme: Zeilen mit \textbf{(E)} gibt es nur für das richtige Ergebnis, nie als Folgepunkt. Die Zeilen «Typische Fehler» sagen, welche Rasterzeile bei diesem Fehler 0 Punkte gibt — sie sind kein zusätzlicher Abzug.\par
3. Andere richtige Wege zählen voll, zum Beispiel beim Grundstück Zerlegen in Rechteck und zwei Dreiecke statt der Trapezformel, beim Regal die Fläche des kleinen Dreiecks direkt aus seinen Katheten ($30$ und $20$) statt über $k^2$, beim gleichschenkligen Dreieck $15 \cdot 36$ statt $\tfrac{30 \cdot 36}{2}$. Gleichwertig: $\sqrt{900} = 30$ und $30^2 = 900$ als Probe; Brüche und Dezimalzahlen mit gleichem Wert.\par
4. Begründe jeden Abzug in einem Satz und nenne die Fehlerart (zum Beispiel «schräge Seite statt Höhe», «halbieren vergessen», «Kathete mit plus gerechnet», «ganze statt halbe Basis», «Wurzel vergessen», «Flächenfaktor k statt k²»). Schreib die Musterlösung nicht ab — zeig mir, wo mein Fehler liegt.\par
5. Gib zum Schluss eine Tabelle «Aufgabe | Punkte | Begründung», die Gesamtpunktzahl (von 25) und — nach der Selbsteinschätzung in Abschnitt 4 — welches Kapitel des Leitprogramms ich wiederholen soll. Keine Note.\par
6. Fehlt ein Bild oder ist eine Aufgabe unlesbar, sag es, statt sie zu bewerten.
\end{tcolorbox}

\needspace{16\baselineskip}
\section*{3 · Musterlösung und Punkteraster}
\textbf{Allgemein:} Ganze Punkte je Zeile. Ein Fehler wird nur einmal abgezogen (Folgefehler). Gleichwertige Wege zählen voll.
In der Spalte P: \textbf{(E)}~= Ergebnispunkt, Ergebnis richtig \emph{und} Weg erkennbar, direkt aus den Angaben; ohne Zeichen = die
Zeile bewertet den Weg, die Skizze oder die Begründung. Der ganze Test ist ohne Taschenrechner gestellt, alle Wurzeln gehen auf.
\textbf{Einheiten:} Fehlt nur die Einheit, gibt es keinen Abzug. \textbf{Ansatzzeilen} bewerten die richtige Gleichung; ein Rechenfehler danach kostet nur die Ergebniszeile.
«Typische Fehler» nennt, welche Zeile 0 Punkte gibt, und die Punkte, die bleiben.
\textbf{Teilrichtige Zeilen:} Verlangt eine Zeile mehrere Angaben und stimmt nur ein Teil, gibt die Zeile 0 Punkte,
ausser die Zeile nennt ausdrücklich eine Aufteilung.

\begin{lpaufgabe}{G1}{Grundstück}{3 P}
\begin{raster}
Formel & 1 & Trapez: $A = \dfrac{a + c}{2} \cdot h$ mit den parallelen Seiten $15$ und $9$ und der Höhe $8$ (nicht $8.2$ oder $8.9$)\\
Rechnung & 1 & $\dfrac{15 + 9}{2} \cdot 8 = 12 \cdot 8$\\
Ergebnis & 1\E & $96$~m$^2$ (auch zerlegt: $9 \cdot 8 + \tfrac{2 \cdot 8}{2} + \tfrac{4 \cdot 8}{2} = 72 + 8 + 16 = 96$)\\
\end{raster}
\fehler{Schräge Seite als Höhe: Zeile Formel 0, Rechnung als Folgepunkt $\to$ 1 von 3\,P. $(15 + 9) \cdot 8 = 192$ (nicht halbiert): Zeilen Rechnung und Ergebnis 0 $\to$ 1 von 3\,P.}
\end{lpaufgabe}

\begin{lpaufgabe}{G2}{Dreieck und Parallelogramm}{3 P}
\begin{raster}
(a) Ansatz & 1 & $\dfrac{g \cdot 7}{2} = 42$, also $7g = 84$\\
(a) Ergebnis & 1\E & $g = 12$~cm\\
(b) & 1\E & $9 \cdot 6 = 54$~cm$^2$; die schräge Seite $7.5$~cm wird nicht gebraucht\\
\end{raster}
\fehler{(a) $42 : 7 = 6$ (Halbieren vergessen): Zeilen (a) 0 $\to$ 1 von 3\,P. (b) $9 \cdot 7.5 = 67.5$ (schräge Seite): Zeile (b) 0 $\to$ 2 von 3\,P.}
\end{lpaufgabe}

\begin{lpaufgabe}{G3}{Pythagoras rechnen und prüfen}{3 P}
\begin{raster}
(a) Ansatz & 1 & $c^2 = 18^2 + 24^2 = 324 + 576 = 900$\\
(a) Ergebnis & 1\E & $c = \sqrt{900} = 30$~cm\\
(b) & 1 & $8^2 + 15^2 = 64 + 225 = 289$, aber $18^2 = 324$: nicht rechtwinklig (Rechnung und Schluss verlangt)\\
\end{raster}
\fehler{(a) $18 + 24 = 42$ (Seiten addiert): Zeilen (a) 0 $\to$ 1 von 3\,P. (a) Wurzel vergessen ($c = 900$): Zeile (a) Ergebnis 0 $\to$ 2 von 3\,P.
(b) «ja, weil $18$ die längste Seite ist» ohne Rechnung: Zeile (b) 0 $\to$ 2 von 3\,P.}
\end{lpaufgabe}

\begin{lpaufgabe}{G4}{Hypotenuse erkennen}{3 P}
\begin{raster}
Hypotenuse & 1 & $61$~cm, denn sie liegt gegenüber dem rechten Winkel (zwischen $11$ und $x$)\\
Ansatz & 1 & $x^2 = 61^2 - 11^2 = 3721 - 121 = 3600$\\
Ergebnis & 1\E & $x = 60$~cm\\
\end{raster}
\fehler{$x$ als Hypotenuse ($x^2 = 61^2 + 11^2 = 3842$): Zeilen Hypotenuse, Ansatz und Ergebnis 0 $\to$ 0 von 3\,P. $61 - 11 = 50$: Zeilen Ansatz und Ergebnis 0 $\to$ 1 von 3\,P.}
\end{lpaufgabe}

\begin{lpaufgabe}{G5}{Gleichschenkliges Dreieck}{4 P}
\begin{raster}
(a) Skizze & 1 & Höhe von der Spitze senkrecht auf die Basis; Hilfsdreieck mit halber Basis $15$~cm, Höhe $h$ und Schenkel $39$~cm als Hypotenuse\\
(b) Ansatz & 1 & $h^2 = 39^2 - 15^2 = 1521 - 225 = 1296$\\
(b) Ergebnis & 1\E & $h = 36$~cm\\
(c) & 1 & $\dfrac{30 \cdot 36}{2} = 540$~cm$^2$ (oder $15 \cdot 36$). Auch als Folgepunkt mit der eigenen Höhe.\\
\end{raster}
\fehler{Ganze Basis $30$ ins Hilfsdreieck ($h^2 = 1521 - 900 = 621$, keine ganze Zahl): Zeilen (a) und (b) Ansatz 0, (b) Ergebnis 0, (c) als Folgepunkt nur mit eigener Höhe $\to$ höchstens 1 von 4\,P.
(c) nicht halbiert ($1080$): Zeile (c) 0 $\to$ 3 von 4\,P.}
\end{lpaufgabe}

\begin{lpaufgabe}{G6}{Wasser im Rohr}{3 P}
\begin{raster}
Abstand & 1 & Radius $25$~cm, Abstand des Wasserspiegels vom Mittelpunkt $d = 25 - 10 = 15$~cm\\
Hilfsdreieck & 1 & Radius als Hypotenuse, $d$ und halbe Breite als Katheten: $\sqrt{25^2 - 15^2} = \sqrt{400} = 20$~cm. Auch als Folgepunkt mit dem eigenen Abstand, solange $0 < d < 25$.\\
Breite & 1 & $2 \cdot 20 = 40$~cm. Auch als Folgepunkt aus der eigenen halben Breite.\\
\end{raster}
\fehler{Durchmesser $50$ als Hypotenuse: Zeile Hilfsdreieck 0, Breite als Folgepunkt $\to$ 2 von 3\,P. Halbe Breite nicht verdoppelt ($20$~cm): Zeile Breite 0 $\to$ 2 von 3\,P.
$d = 10$ (Tiefe als Abstand genommen): Zeile Abstand 0, Folgepunkte mit $d = 10$ ($2\sqrt{525}$) nur, wenn richtig angesetzt $\to$ höchstens 2 von 3\,P.}
\end{lpaufgabe}

\begin{lpaufgabe}{G7}{Eckregal}{3 P}
\begin{raster}
(a) & 1\E & $k = \dfrac{30}{90} = \dfrac{1}{3}$, kürzere Kathete $60 \cdot \tfrac{1}{3} = 20$~cm (beides verlangt)\\
(b) & 1 & $\dfrac{90 \cdot 60}{2} \cdot \left(\tfrac{1}{3}\right)^2 = 2700 \cdot \tfrac{1}{9} = 300$~cm$^2$ (oder $\tfrac{30 \cdot 20}{2} = 300$)\\
(c) & 1 & $2700 - 300 = 2400$~cm$^2$. Auch als Folgepunkt aus den eigenen Flächen.\\
\end{raster}
\fehler{(b) $2700 \cdot \tfrac{1}{3} = 900$ (Flächenfaktor $k$ statt $k^2$): Zeile (b) 0, (c) als Folgepunkt ($1800$) $\to$ 2 von 3\,P.
$k = \tfrac{30}{60} = \tfrac{1}{2}$ (falsche Seiten zugeordnet): Zeile (a) 0, (b) und (c) als Folgepunkte ($675$, $2025$) $\to$ 2 von 3\,P.
$k = 3$ (Kehrwert): ausgesägte Fläche grösser als das Regal, unmöglich: keine Folgepunkte $\to$ 0 von 3\,P.}
\end{lpaufgabe}

\begin{lpaufgabe}{G8}{Schatten}{3 P}
\begin{raster}
(a) Skizze & 1 & zwei rechtwinklige Dreiecke (Stab mit Schatten, Turm mit Schatten), Sonnenstrahlen parallel; Stab entspricht Turm, Stabschatten entspricht Turmschatten\\
(b) Ansatz & 1 & $\dfrac{h}{14} = \dfrac{1.2}{0.8}$ oder $k = \dfrac{14}{0.8} = 17.5$\\
(b) Ergebnis & 1\E & $h = 1.5 \cdot 14 = 21$~m\\
\end{raster}
\fehler{Unterschied statt Verhältnis ($14 + 0.4 = 14.4$~m): Zeilen (b) 0 $\to$ 1 von 3\,P. Verhältnis umgekehrt ($14 \cdot \tfrac{0.8}{1.2} \approx 9.3$~m): Zeile (b) Ansatz 0, Ergebnis 0 $\to$ 1 von 3\,P.}
\end{lpaufgabe}

\needspace{12\baselineskip}
\section*{4 · Selbsteinschätzung}
\begin{tabularx}{\linewidth}{@{}l X@{}}
\toprule
\textbf{22 – 25 P} & Die geprüften Teile sitzen. Wo du Punkte verloren hast: das Kapitel dieser Aufgabe nochmals.\\
\textbf{17 – 21 P} & Den schwächsten Teil nochmals, gemessen am Anteil (Teil A 6, B 6, C 7, D 6 P): Tüfteln und Übungen des Kapitels dieses Teils.\\
\textbf{11 – 16 P} & Zurück zu den Kapiteln aller Aufgaben, in denen du Punkte verloren hast.\\
\textbf{0 – 10 P} & Zurück zu den Kapiteln jedes Teils, in dem du weniger als die Hälfte erreicht hast; sind es alle vier, zum Vortest und zu Kapitel 1.\\
\bottomrule
\end{tabularx}

\smallskip
\textbf{Aufgabe $\to$ Kapitel:} G1, G2 $\to$ 1; G3, G4 $\to$ 2; G5, G6 $\to$ 3; G7, G8 $\to$ 4.
\end{document}
